\chapter{متغیرها و انواعِ داده}%
\label{ch:variables}

در فصلِ \ref{ch:algorithms}، هنگامِ طراحیِ الگوریتمِ «بزرگ‌ترین عدد»، به
جایی برای نگهداریِ «بزرگ‌ترینِ فعلی» نیاز پیدا کردیم. برنامه‌ها تقریباً
همیشه باید چیزهایی را به خاطر بسپارند: امتیازِ بازیکن، نامِ کاربر، مجموعِ
نمره‌ها. ابزارِ این «به خاطر سپردن»، \textbf{متغیر} است.

\section{متغیر؛ جعبه‌ای با نام و برچسب}

متغیر را مانندِ جعبه‌ای تصور کنید که رویش نامی نوشته‌اید و درونش یک مقدار
گذاشته‌اید. هر وقت آن نام را در برنامه بیاورید، یعنی «محتوای این جعبه».
ساختنِ متغیر در سلام چنین است:

\begin{salam}
تابع اصلی:
    متغیر امتیاز : صحیح = 100
    چاپ امتیاز
تمام
\end{salam}

\begin{progout}
100
\end{progout}

خطِ دوم را چنین بخوانید: «\textbf{متغیر}ی به نامِ \textbf{امتیاز} بساز که
نوعش \textbf{صحیح} (عددِ درست) است و مقدارِ \textbf{۱۰۰} را در آن بگذار.»
ساختار همیشه یکسان است:

\begin{center}
\code{متغیر} \;\textit{نام}\;:\;\textit{نوع} \;=\; \textit{مقدار}
\end{center}

دقت کنید که در \code{چاپ امتیاز} خبری از نقل‌قول نیست: می‌خواهیم
\emph{محتوای} جعبه چاپ شود (۱۰۰)، نه واژهٔ «امتیاز».

\section{تغییر دادنِ مقدار}

اسمش رویش است: متغیر می‌تواند تغییر کند. با نوشتنِ \code{نام = مقدارِ تازه}
محتوای جعبه عوض می‌شود (این بار دیگر واژهٔ \code{متغیر} را نمی‌نویسیم، چون
جعبه از قبل ساخته شده):

\begin{salam}
تابع اصلی:
    متغیر امتیاز : صحیح = 100
    چاپ "امتیاز اول:", امتیاز
    امتیاز = 250
    چاپ "امتیاز بعدی:", امتیاز
    امتیاز = امتیاز + 50
    چاپ "با جایزه:", امتیاز
تمام
\end{salam}

\begin{progoutfa}
امتیاز اول: 100
امتیاز بعدی: 250
با جایزه: 300
\end{progoutfa}

خطِ \code{امتیاز = امتیاز + 50} برای تازه‌کارها عجیب است: مگر می‌شود چیزی
با خودش به‌علاوهٔ ۵۰ برابر باشد؟! راز اینجاست که علامتِ \code{=} در
برنامه‌نویسی «مساوی بودن» نیست؛ \textbf{انتساب} است، یعنی: «اول سمتِ راست
را حساب کن، بعد نتیجه را در جعبهٔ سمتِ چپ بگذار.» پس سمتِ راست با مقدارِ
فعلی (۲۵۰) حساب می‌شود و می‌شود ۳۰۰، بعد ۳۰۰ در جعبه گذاشته می‌شود. این
همان نمادِ $\leftarrow$ شبه‌کدهای فصلِ \ref{ch:algorithms} است.

\begin{note}
ترتیبِ خط‌ها اهمیت دارد. برنامه از بالا به پایین اجرا می‌شود و هر
\code{چاپ} مقدارِ \emph{همان لحظهٔ} متغیر را نشان می‌دهد؛ به همین دلیل سه
خروجیِ متفاوت دیدیم.
\end{note}

\section{ثابت: جعبه‌ای که قفل است}

بعضی مقدارها نباید هرگز عوض شوند: تعدادِ روزهای هفته، عددِ پی، سالِ تولد.
برای این‌ها به‌جای \code{متغیر} از \code{ثابت} استفاده کنید:

\begin{salam}
تابع اصلی:
    ثابت روزهای هفته : صحیح = 7
    چاپ روزهای هفته
تمام
\end{salam}

اگر جایی از برنامه سهواً بنویسید \code{روزهای هفته = 8}، کامپایلر خطا
می‌گیرد و جلوی اشتباه را همان اول می‌گیرد. راهِ سومی هم هست: اگر هیچ
واژه‌ای ننویسید و مستقیم با نام شروع کنید (\code{قد : اعشار = 1.75})،
باز هم یک مقدارِ تغییرناپذیر ساخته‌اید. قاعدهٔ خوب این است: \emph{مگر واقعاً
به تغییر نیاز دارید، تغییرناپذیر بسازید}؛ برنامه‌ای که جعبه‌های کمتری در آن
عوض می‌شوند، کم‌اشتباه‌تر است.

\section{انواعِ داده}

تا اینجا فقط با عددِ صحیح کار کردیم، اما داده‌ها جورواجورند. مهم‌ترین
انواعِ سلام این‌هاست:

\begin{center}
\begin{tabular}{lll}
\toprule
نوع & چه چیزی نگه می‌دارد & نمونه \\
\midrule
\code{صحیح}  & عددِ درست (بی‌اعشار)         & \code{42} یا \code{-7} \\
\code{اعشار} & عددِ اعشاری                  & \code{3.14} یا \code{-0.5} \\
\code{رشته}  & متن                          & \code{"سلام"} \\
\code{منطقی} & فقط \code{درست} یا \code{نادرست} & \code{درست} \\
\code{نویسه} & یک حرفِ تنها                 & \code{'م'} \\
\bottomrule
\end{tabular}
\end{center}

\begin{salam}
تابع اصلی:
    نام : رشته = "سارا"
    متغیر سن : صحیح = 12
    قد : اعشار = 1.48
    عضو باشگاه : منطقی = درست
    چاپ نام, سن, قد, عضو باشگاه
تمام
\end{salam}

\begin{progoutfa}
سارا 12 1.48 true
\end{progoutfa}

نوعِ \code{منطقی} شاید غریب به نظر برسد اما ستارهٔ فصلِ
\ref{ch:conditionals} است: هر جا برنامه بخواهد «تصمیم» بگیرد، پای یک مقدارِ
منطقی در میان است.

\subsection{چرا نوع مهم است؟}

نوع به کامپایلر می‌گوید با محتوای جعبه چه کارهایی مجاز است. دو عددِ صحیح را
می‌شود ضرب کرد، اما ضربِ دو رشته بی‌معنی است. اگر بنویسید
\code{"سلام" * 3}، کامپایلر همان‌جا خطا می‌گیرد. این سخت‌گیری دوستِ شماست:
بسیاری از اشتباه‌ها پیش از اجرا، همان لحظهٔ ساختنِ برنامه، پیدا می‌شوند.

\subsection{خودکار: بگذارید سلام حدس بزند}

وقتی مقدارِ اولیه جلوی چشم است، نوشتنِ نوع گاهی تکراری است؛ معلوم است که
\code{"سارا"} رشته است. با نوعِ \code{خودکار} تشخیص را به سلام می‌سپارید:

\begin{salam}
تابع اصلی:
    نام : خودکار = "سارا"
    سن : خودکار = 12
    چاپ نام, سن
تمام
\end{salam}

جعبه همچنان نوعِ دقیقِ خودش را دارد (رشته، صحیح)؛ فقط زحمتِ نوشتنش را سلام
کشیده است.

\section{نام‌گذاریِ خوب}

نامِ متغیر را شما انتخاب می‌کنید و این انتخاب مهم‌تر از آن است که به نظر
می‌رسد. برنامه‌ای با نام‌های \code{الف} و \code{ب} و \code{ج} را یک هفته
بعد خودتان هم نمی‌فهمید. چند قاعدهٔ سرانگشتی:

\begin{itemize}
  \item نام باید بگوید جعبه \emph{چه چیزی} نگه می‌دارد: \code{مجموع نمرات}
        بهتر از \code{م}، و \code{تعداد دانش آموزان} بهتر از \code{ت} است.
  \item نام با عدد شروع نمی‌شود و علامت‌هایی مانندِ \code{+} و \code{-} در
        آن راه ندارند.
  \item دو متغیرِ هم‌نام نمی‌توان ساخت؛ هر جعبه برچسبِ یکتای خود را دارد.
\end{itemize}

\begin{deep}[نگاهی به درون: جعبه‌ها کجا هستند؟]
جعبه‌ها در \emph{حافظهٔ اصلیِ} رایانه (\lr{RAM}) ساخته می‌شوند. حافظه را
مانندِ دیواری از میلیون‌ها خانهٔ کوچکِ شماره‌دار تصور کنید. وقتی می‌نویسید
\code{متغیر امتیاز : صحیح = 100}، سلام چند خانهٔ کنارِ هم را برمی‌دارد،
عددِ ۱۰۰ را در آن‌ها می‌گذارد و نامِ «امتیاز» را به نشانیِ آن خانه‌ها گره
می‌زند. نوعِ متغیر تعیین می‌کند چند خانه لازم است و الگوی صفر و یک‌ها چگونه
تفسیر شود.
\end{deep}

\section{تمرین‌ها}

\exercise سه متغیر بسازید: نامِ کتابِ محبوبتان (رشته)، تعدادِ صفحه‌هایش
(صحیح) و امتیازی که به آن می‌دهید (اعشار، از ۵). هر سه را در یک خط چاپ
کنید.

\exercise متغیری به نامِ \code{موجودی} با مقدارِ ۵۰۰ بسازید. با سه خطِ
جداگانه: ۲۰۰ به آن اضافه کنید، بعد ۱۵۰ از آن کم کنید و در پایان چاپش کنید.
پیش از اجرا حدس بزنید چه عددی چاپ می‌شود.

\exercise این برنامه را بدونِ اجرا دنبال کنید و بگویید چه چاپ می‌کند:
\begin{salam}
تابع اصلی:
    متغیر آ : صحیح = 4
    متغیر ب : صحیح = 9
    آ = ب
    ب = آ
    چاپ آ, ب
تمام
\end{salam}
اگر هدف «جابه‌جا کردنِ» مقدارِ دو جعبه بود، این برنامه موفق شده است؟ چرا؟

\exercise با یک متغیرِ \emph{سوم} برنامهٔ تمرینِ قبل را طوری بازنویسی کنید
که مقدارها واقعاً جابه‌جا شوند (راهنمایی: قبل از خراب کردنِ \code{آ}،
مقدارش را جایی امن نگه دارید).

\exercise کدام‌یک برای «تعدادِ خواهر و برادرها» نوعِ مناسب‌تری است: صحیح یا
اعشار؟ برای «وزن» چطور؟ برای «آیا چراغ روشن است؟» چطور؟
